\documentclass[11pt]{article}
\usepackage[margin=1in]{geometry}
\usepackage{amsmath,amssymb}
\usepackage[hidelinks]{hyperref}

\title{Literature Status Note:\\Open Problem A3 in arXiv:2001.02226}
\author{fa\_banach\_001 lane 0}
\date{2026-07-19}

\begin{document}
\maketitle

\section*{Status}

\textbf{literature\_already\_answered}.  This packet records an exact later
literature answer, not a new proof produced by the run.

\section*{Source Question}

The source is Petru Mironescu and Jean Van Schaftingen,
``Trace theory for Sobolev mappings into a manifold'', arXiv:2001.02226,
published in \emph{Ann. Fac. Sci. Toulouse Math. (6)} 30 (2021), 281--299.
In the introduction, Open Problem A3 asks whether the local trace space
\[
  W^{1-1/p,p}(B^{m-1},N)
\]
has the extension property when \(4 \le p < m\),
\[
  \pi_1(N),\ldots,\pi_{\lfloor p-2\rfloor}(N)
  \quad\text{are finite,}
  \qquad
  \pi_{\lfloor p-1\rfloor}(N)=0.
\]
The paper immediately notes that Bethuel--Demengel had conjectured a positive
answer.

\section*{Answering Source}

The supporting paper is Jean Van Schaftingen, ``The extension of traces for
Sobolev mappings between manifolds'', arXiv:2403.18738.  Its introduction says
that the paper gives a complete answer to trace surjectivity and proves that
there is no further obstruction beyond the known topological and analytical
ones.  The same paper cites the 2021 Mironescu--Van Schaftingen paper.

The decisive local theorem is Theorem 1.1, labelled
\texttt{theorem\_extension\_halfspace}.  For \(1<p<m\), it states that
\[
  \operatorname{tr}_{\mathbb R^{m-1}}
  \bigl(\dot W^{1,p}(\mathbb R^m_+,N)\bigr)
  =
  \dot W^{1-1/p,p}(\mathbb R^{m-1},N)
\]
if and only if
\[
  \pi_1(N),\ldots,\pi_{\lfloor p-2\rfloor}(N)
  \quad\text{are finite,}
  \qquad
  \pi_{\lfloor p-1\rfloor}(N)=0.
\]
Moreover the theorem supplies a constant \(C\) so that the extension can be
chosen with
\[
  \int_{\mathbb R^m_+} |DU|^p
  \le
  C \iint_{\mathbb R^{m-1}\times\mathbb R^{m-1}}
  \frac{d(u(y),u(z))^p}{|y-z|^{p+m-2}}\,dy\,dz .
\]
The collar formulation, Theorem 1.1', gives the corresponding local manifold
version.

\section*{Identification}

Open Problem A3 is exactly the subrange \(4\le p<m\) of the local
surjectivity theorem above.  Its homotopy hypotheses match the necessary and
sufficient hypotheses in the 2024 theorem.  Therefore A3 has a positive
literature answer.  The adjacent quantitative question in the source paper is
also answered positively in the half-space/collar setting by the displayed
linear estimate.

The source paper formulates the question on \(B^{m-1}\) after explaining that
trace theory has local versions in which \(\mathbb R^{m-1}\) is replaced by a
Lipschitz domain.  The 2024 half-space and collar theorems are the standard
local model for that same trace-extension problem.

\section*{Scope and Limitations}

This packet does not claim a new solution.  It records that the target should
be skipped by future agents as already answered in later literature.  The
global compact-manifold theorem in arXiv:2403.18738 has an additional
skeleton-extension condition involving \(M\) and \(\partial M\); this note only
identifies the local A3/linear-estimate answer.

\begin{thebibliography}{9}

\bibitem{MVS2021}
P. Mironescu and J. Van Schaftingen,
\emph{Trace theory for Sobolev mappings into a manifold},
Ann. Fac. Sci. Toulouse Math. (6) 30 (2021), 281--299.
arXiv:2001.02226.

\bibitem{VS2024}
J. Van Schaftingen,
\emph{The extension of traces for Sobolev mappings between manifolds},
arXiv:2403.18738.

\end{thebibliography}

\end{document}

